\chapter{INTRODUCTION}

\section{Background and Research Motivation}

University students face a combination of academic pressure, social comparison, family expectations, career uncertainty, loneliness, and burnout during a period when many mental-health difficulties first become visible. Large-scale studies of college students show that mental-health concerns are common in this population \citep{Auerbach2018CollegeMentalHealth}, while recent trend analyses report continuing concerns around student mental health and help-seeking \citep{Lipson2022CollegeTrends}. At the same time, students may delay or avoid formal counseling because of stigma, limited access, cost, time constraints, or uncertainty about whether their distress is serious enough to discuss with a professional. Stigma and help-seeking barriers among college students have been documented empirically \citep{Eisenberg2009StigmaHelpSeeking}.

Digital support tools are therefore attractive because they can be low-barrier, always available, and easier to approach than formal services. Early conversational-agent work such as Woebot explored automated CBT-oriented support for young adults with symptoms of depression and anxiety \citep{Fitzpatrick2017Woebot}. More recent LLM-based systems such as SoulChat show that language models can be adapted toward empathetic listening, comfort, and client-centered dialogue \citep{Chen2023SoulChat}. These systems suggest that conversational AI can provide useful supportive interaction, especially when users need reflection, organization of thoughts, or encouragement to seek stronger real-world support.

However, mental-health dialogue is not simply a problem of producing warm text. A response can be fluent, empathetic, and still clinically poorly timed. For example, a user who is physiologically overwhelmed may need grounding before cognitive analysis; a user who has already identified an automatic thought may need guided discovery rather than another broad request to describe feelings. Work on cognitive reframing and CBT-oriented counseling conversations suggests that therapeutic technique and response timing matter \citep{Sharma2023CognitiveReframing,Lee2024Cactus}. This thesis is motivated by the gap between fluent chatbot conversation and structured, inspectable, safety-aware supportive dialogue.

\section{Research Problem and Gap}

Many LLM-based support systems follow a one-to-one design: the user sends a message, and a single model produces a response. This design is simple, but it compresses many responsibilities into one component. The same model is expected to understand the user's state, choose a therapeutic direction, maintain context, generate the response, and respect safety boundaries. In student mental-health conversations, this can lead to generic reassurance, repeated validation, premature advice, route mixing, or poor stage progression.

The first gap is a process-control gap. Existing LLM counseling systems can generate supportive responses, but it remains difficult to ensure that the conversation follows a coherent therapeutic route across multiple turns. A student-support system should be able to hold, advance, or regress a stage based on evidence in the conversation. It should also avoid mixing CBT, mindfulness-based intervention, and behavioral activation when the user's current need points to only one of them. Behavioral assessment work argues that therapist-like systems should be evaluated through observable counseling behavior rather than general chatbot quality alone \citep{Chiu2024LLMTherapists}.

The second gap is a virtual-group gap. Human group therapy theory emphasizes factors such as universality, catharsis, hope, and interpersonal learning \citep{Yalom2020GroupPsychotherapy}. Yet most LLM mental-health systems remain one-to-one chatbots. Adding multiple agents without control would not solve this problem; it could make the dialogue repetitive, confusing, or unsafe. The open question is how to operationalize peer-like support so that virtual peers can contribute when useful while the therapist remains the final responsible voice.

The third gap is an evaluation gap. If an evaluation only rewards fluency, empathy, or general helpfulness, it may miss whether the system followed the intended modality, advanced the conversation appropriately, used peer support at the right time, or handled safety boundaries. LLM-as-a-judge methods can support scalable evaluation, but they also require caution and calibration \citep{Zheng2023LLMJudge}. Mental-health counseling and psychosocial safety work further emphasizes the need to test boundary control, robustness, and sensitive-risk scenarios \citep{Li2025CounselBench,Luo2025DialogGuard}. This thesis therefore treats system design and benchmark design as connected problems.

\section{Research Objectives, Questions, and Scope}

The objective of this thesis is to design and evaluate a controllable AI support prototype for Vietnamese student mental-health conversations. The work focuses on dialogue structure, agent coordination, inspectability, and safety boundaries. It does not attempt to prove treatment outcomes, replace professional counseling, or deploy a clinical product.

The thesis is guided by three research questions:

\begin{enumerate}
    \item \textbf{RQ1:} How can an LLM-based student-support system maintain route and stage control across CBT, mindfulness-based intervention, and behavioral activation?
    \item \textbf{RQ2:} How can Yalom-inspired peer agents contribute to a virtual group-support experience without taking clinical authority away from the therapist?
    \item \textbf{RQ3:} How can such a system be evaluated beyond response fluency, using safety, modality fidelity, group dynamics, context retention, and multi-turn progress?
\end{enumerate}

To answer these questions, the thesis proposes a blackboard-based multi-agent virtual group therapy prototype. The conversation begins with onboarding and screening. DASS-21 is used as a practical screening signal for initial routing, supported by its original formulation and Vietnamese youth validation work \citep{Lovibond1995DASS,Le2017DASS21Vietnamese}. The system then routes the conversation into CBT, mindfulness-based intervention, or behavioral activation. CBT is included for cognitive restructuring and guided discovery \citep{Hofmann2012CBTMetaAnalyses}; mindfulness-based intervention is included for grounding and decentering in high-arousal or ruminative states \citep{Zenner2014MindfulnessSchools}; and behavioral activation is included for low-energy, avoidance, and academic-burnout patterns \citep{Mazzucchelli2010BehavioralActivation}.

The system uses a shared blackboard state to coordinate specialized components. A clinical assessor updates route, stage, safety flags, and case formulation. Two peer agents, Nam and Linh, observe the blackboard and may draft short supportive contributions only when the current stage and Yalom factor allow it. Nam is designed for universality and catharsis; Linh is designed for hope and interpersonal learning. The therapist orchestrator remains the final decision maker: it reviews peer drafts, rewrites or discards them when needed, and produces the user-visible response. This design follows multi-agent and blackboard coordination principles while avoiding uncontrolled multi-speaker dialogue \citep{Wu2023AutoGen,Salemi2025BlackboardSystem}.

The thesis makes four main contributions:

\begin{enumerate}
    \item \textbf{A stage-constrained dialogue architecture} for CBT, mindfulness-based intervention, and behavioral activation, including route goals, transition logic, and response validation.
    \item \textbf{A therapist-governed blackboard coordination mechanism} that separates assessment, peer observation, therapist planning, guardrails, and memory update through shared state.
    \item \textbf{A Yalom-informed peer policy} in which Nam and Linh provide limited, therapist-approved support while preserving silence when peer contribution would harm the therapeutic flow.
    \item \textbf{A VSM benchmark workflow} for multi-turn evaluation of safety, modality fidelity, group dynamics, conversation progress, reliability, confidence intervals, and human-audit readiness.
\end{enumerate}

The scope of the thesis is deliberately limited. The system is a research prototype for supportive dialogue, not a diagnostic instrument or a substitute for licensed care. The benchmark is designed for engineering evaluation and does not establish treatment outcomes. Real-world use would require clinician review, stronger baseline comparison, privacy analysis, safety monitoring, and formal ethical evaluation.

\section{Thesis Structure}

The remainder of this thesis is organized as follows. Chapter 2 reviews screening measures, therapeutic foundations, LLM-based mental-health systems, Yalom-inspired group support, blackboard multi-agent coordination, and evaluation methods. It also identifies the process-control, virtual-group, and evaluation gaps addressed by this work.

Chapter 3 presents the methodology and system design. It describes the blackboard state, dialogue graph, clinical assessor, stage contracts, case formulation layer, peer agents, therapist orchestrator, guardrails, memory updater, and implementation interfaces.

Chapter 4 describes the experimental setup and evaluation framework. It introduces the VSM benchmark, dataset splits, compared systems, metrics, hybrid scoring, structural results, and staged evaluation plan.

Chapter 5 concludes the thesis by summarizing the contributions, discussing limitations and ethical considerations, and outlining future work for stronger evaluation and safer deployment.
